%
% display_linkage.tex
%
% A Z specification of the corrected display-linkage design: the Hub holds
% content it cannot deliver, retries at a bounded and eventually slow
% cadence, and never touches the display's process lifecycle.  Two
% concurrent Hub-side actors dial the same client registry -- the
% HubReplicator (the content sender) and DisplayLiveness (the cheap,
% frequent prober) -- coordinated by one shared reconnect signal.
%
% Models the DESIGN in
% docs/architecture/display-presence-demand-driven.md, sections 4
% (DisplayLinkage), 6.1 (the three-way cycle outcome), 6.2 (the two
% independently-tuned backoff curves), 6.3 (the interruptible wait,
% coordinated through ClientRegistry._reconnected and NOT holding
% ClientRegistry._lock across it), and 6.4 (DisplayLiveness's two-speed
% pacing).  It does not re-model the single-writer send/reap/reconnect/menu
% discipline that docs/hub_replicator.tex already proves for one worker
% against one display; see the note below on why this is a separate
% specification rather than an extension of that one.
%
% Type-check:  fuzz docs/display_linkage.tex
% Model-check (see the Verification section for the exact goal predicates):
%   probcli docs/display_linkage.tex -model_check -p DEFAULT_SETSIZE 2
%
% Formatting follows the fuzz house rules established by hub_replicator.tex:
% \quad~ for continuation indent, schema lines under 80 columns, predicates
% broken at \land/\lor/\implies.
%

\documentclass[a4paper,10pt,fleqn]{article}
\usepackage[margin=1in]{geometry}
\usepackage{fuzz}
\usepackage[colorlinks=true,linkcolor=blue,citecolor=blue,urlcolor=blue]{hyperref}

\begin{document}

\title{Display Linkage: a Z Specification}
\author{Formal model of the Hub's hold/reconcile/never-control discipline}
\date{September 2026}
\maketitle

% ============================================================
\section{Introduction}
% ============================================================

Closing the Lux window used to black-hole content: an idle, no-display Hub
retried at a fixed, tight cadence forever, because the code could not tell
``no display has ever been connected'' apart from ``a connected display just
misbehaved,'' and both conditions advanced the same short backoff.  The
corrected design (\texttt{docs/architecture/display-presence-demand-driven.md})
holds content safely in \texttt{HubDisplay} (unchanged), reconciles it the
moment a display connects (unchanged, DES-068), and fixes the one thing that
was actually broken: the retry \emph{cadence} when nothing is connected at
all, and the coordination between the Hub's two independent threads that
each dial the display.

\paragraph{Why a new specification, not an extension of \texttt{hub\_replicator.tex}.}
\texttt{docs/hub\_replicator.tex} already proves six invariants for the
Hub's \emph{single} background worker sending to \emph{one} display: no
second writer, agent-path liveness, no lock cycle between the store lock and
the client lock, no lost update across a reap or reconnect, no two live
displays, and menu durability.  That model has exactly one thread that ever
touches the display connection.  This design's entire novelty is a
\emph{second} thread --- \texttt{DisplayLiveness} --- dialing the
\emph{same} client registry concurrently, coordinated with the replicator
through a shared \texttt{threading.Event} that is deliberately consulted
\emph{without} holding the registry's lock, plus two independently-tuned
exponential backoff curves that must never be conflated.  None of that
concurrency shape, nor any of its invariants, exists in
\texttt{hub\_replicator.tex}'s \texttt{Repl} schema.  Grafting a second actor
onto an already-verified, committed regression artifact would force
re-deriving its six proven invariants under an interleaving they were never
checked against, for invariants that are otherwise orthogonal to them.  A
fresh, narrower specification --- cross-referencing
\texttt{hub\_replicator.tex} for the reconciliation hook it already models
(\texttt{Ensure}/\texttt{Remark}/\texttt{Reconn}'s re-mark, the concrete
counterpart of this design's DES-068 connect-success reconcile) --- is more
faithful to what changed and leaves the existing proof untouched.  (The
mission brief cites this as ``\S10 item 13''; the design document's own
numbering makes it item 13's cross-reference to item 11, the test list ---
noted here since the two numbers do not otherwise agree.)

\paragraph{Scope.} This model covers the connect/disconnect/backoff
mechanism newly introduced by the design: \texttt{DisplayLinkage} as an
observed classification of \texttt{(connected, live content)}; the two
concurrent actors dialing \texttt{ClientRegistry}; the disconnected backoff
(2s~$\to$~120s in the design, abstracted here to a small bounded step count)
kept structurally separate from the wedged-display backoff (0.1s~$\to$~2.0s,
likewise abstracted); and the shared reconnect signal with its
lock-freedom property.  It does \emph{not} re-derive the fine-grained
send-cycle mechanics (draining, snapshotting under a store lock, per-scene
blank-on-empty) that \texttt{hub\_replicator.tex} already covers: one
``send cycle'' here is modelled as a single atomic step that delivers
whatever is currently held, which is adequate for every invariant below and
is stated plainly so the simplification is never mistaken for silence.

Eight properties, drawn directly from \S12 of the design document, must hold
across every interleaving of the two actors and the content-owning agents:

\begin{description}
  \item[I1 --- content never lost while disconnected.] Every scene that is
    live is either queued for delivery or already known to be shown; no
    interleaving of push, disconnect, and retry silently drops it.
  \item[I2 --- reconcile exactly once.] The moment either actor's dial
    turns a disconnected registry into a connected one, every currently
    live scene is marked for redelivery in that same step --- never missed,
    never re-queued twice for one reconnect.
  \item[I3 --- neither actor blocks or spins on a missing display.] Every
    dial resolves in one bounded step; the replicator's retry interval is
    always the disconnected backoff's current value; the prober's interval
    is always one of exactly two named speeds.
  \item[I4 --- the disconnected backoff is genuinely exponential and
    reaches a slow steady state.] Under a sustained disconnection its own
    counter climbs to its own, higher cap --- not merely bounded, but
    bounded at the \emph{right}, slower ceiling.
  \item[I5 --- the Hub never issues a process-control call to the
    display.] No Hub-side operation ever causes a display to become
    reachable; only the environment does.
  \item[I6 --- a display connect promptly breaks the replicator out of
    the slow-retry state.] Via the specific mechanism the design defines,
    not merely ``eventually.''
  \item[I7 --- a disconnected cycle never advances the wedged-display
    backoff, and vice versa.] The two counters are structurally disjoint.
  \item[I8 --- deadlock-freedom, with the prober modelled as a genuine
    second actor.] In particular, the specific claim to check: the
    replicator's wait for a reconnect never holds the registry's lock.
\end{description}

% ============================================================
\section{Basic Types}
% ============================================================

\begin{zed}
  [SCENE]
\end{zed}

\texttt{SCENE} is the set of scenes \texttt{HubDisplay} can hold live; one
suffices for every invariant below except I2's ``every currently live
scene'' claim, which needs at least two to show that a batch reconcile
covers more than a singleton.  Bounded by ProB's \texttt{DEFAULT\_SETSIZE}.

% ============================================================
\section{Free Types}
% ============================================================

\begin{zed}
  FLAG ::= set | clear
\end{zed}

\begin{zed}
  LOCK ::= held | free
\end{zed}

Two-valued markers, following \texttt{hub\_replicator.tex}'s own convention:
\texttt{FLAG} for the connected belief, the shared reconnect signal, and the
``woken early or timed out'' witness; \texttt{LOCK} for
\texttt{ClientRegistry.\_lock}.

\begin{zed}
  RPHASE ::= ridle | rwaitdisc | rsending | rhealing
\end{zed}

The replicator's own program counter, restricted to the phases this design
adds or changes.  \texttt{ridle}: no cycle in flight, about to dial.
\texttt{rwaitdisc}: parked in \texttt{\_wait\_disconnected}, having just
failed to connect at all.  \texttt{rsending}: connected, about to deliver
whatever is held.  \texttt{rhealing}: a connected-but-misbehaving delivery
is being reaped/reconnected --- the machinery
\texttt{docs/hub\_replicator.tex} already proves in full; here it is one
atomic step, entered and left.

\begin{zed}
  SPEED ::= pconn | pdisc
\end{zed}

\texttt{DisplayLiveness}'s current probe pacing: \texttt{pconn}, the normal
connected-interaction cadence, or \texttt{pdisc}, the slightly relaxed
\texttt{\_DISCONNECTED\_PROBE\_INTERVAL} --- the two, and only two, named
values \S6.4 defines.

% ============================================================
\section{State}
% ============================================================

The state is flat: one schema, exactly as
\texttt{docs/display\_lifecycle.tex} and \texttt{docs/hub\_replicator.tex}
require for ProB.  Its predicate carries only structural bookkeeping ---
facts true by construction of every operation below, correct or buggy ---
never one of the eight safety properties themselves, which must stay free
to be violated by a faulty variant or the model would prove nothing.

\begin{schema}{Link}
  live : \power SCENE \\
  dirty : \power SCENE \\
  shown : \power SCENE \\
  linked : FLAG \\
  reconnected : FLAG \\
  wokenEarly : FLAG \\
  reachable : FLAG \\
  sendOk : FLAG \\
  regLock : LOCK \\
  rphase : RPHASE \\
  discBackoff : \nat \\
  wedgeBackoff : \nat \\
  lInterval : SPEED
\where
  dirty \subseteq live \\
  shown \subseteq live
\end{schema}

\texttt{live} is the set of scenes currently live in \texttt{HubDisplay} ---
the content-owning agents' own state, touched only by the two content
operations below, never by either dialing actor.  \texttt{dirty} is the set
still owed to the display: exactly the design's ``content is held'' set,
populated by a push and by the reconcile-on-connect, and drained only by an
actual successful delivery.  \texttt{shown} records which live scenes are
believed already delivered --- the abstraction standing in for
\texttt{hub\_replicator.tex}'s own \texttt{displayed} map, coarsened to a
set because this spec does not re-derive per-scene send mechanics.

\texttt{linked} is \texttt{ClientRegistry}'s own belief about whether a
display is connected --- what \texttt{DisplayLinkage.classify} reads as its
\texttt{connected} argument, together with \texttt{\#live} for its
\texttt{live\_scene\_count}.  \texttt{DisplayLinkage} itself is not stored:
consistent with \S4 of the design (``no new stored state, because there is
nothing for the Hub to decide here, only something to observe''), it is a
derived classification of \texttt{(linked, live)}:
\[
  Linkage(linked, live) \defs
  \begin{cases}
    \mathrm{DISCONNECTED} & linked = clear \land live = \emptyset \\
    \mathrm{HELD} & linked = clear \land live \neq \emptyset \\
    \mathrm{CONNECTED\_IDLE} & linked = set \land live = \emptyset \\
    \mathrm{CONNECTED\_ACTIVE} & linked = set \land live \neq \emptyset
  \end{cases}
\]
exactly \texttt{DisplayLinkage.classify}'s own four-way case split.  No
invariant below needs to name a \texttt{Linkage} value directly; each is
stated over \texttt{linked}, \texttt{live}, \texttt{dirty}, and the backoff
state instead, which is equivalent and avoids adding a derived value to the
carrier ProB has to enumerate.

\texttt{reconnected} is the shared \texttt{threading.Event}
(\texttt{ClientRegistry.\_reconnected}): set by whichever actor's dial is
the one that turns \texttt{linked} from \texttt{clear} to \texttt{set},
consumed by the replicator's wait.  \texttt{wokenEarly} records what
\texttt{wait\_for\_reconnect}'s own boolean return value records --- whether
the wait was broken by the event (\texttt{set}) or ran out the clock
(\texttt{clear}) --- the exact fact I6 is stated over.  \texttt{reachable}
and \texttt{sendOk} are the environment's own truth: whether a dial would
succeed right now, and, independently, whether an already-connected send
would succeed right now (a wedged or dead-peer failure is a \emph{send-time}
event, not a dial-time one, so the two are modelled as separate flags rather
than folded into one three-way display-process type).  \texttt{regLock} is
\texttt{ClientRegistry.\_lock}. \texttt{discBackoff} and
\texttt{wedgeBackoff} are the two independently-tuned counters of \S6.2,
abstracted from seconds to a step count; \texttt{lInterval} is
\texttt{DisplayLiveness}'s current pacing.

Modelling \texttt{reachable} and \texttt{sendOk} as two independent flags
rather than one three-valued display-process type (as
\texttt{hub\_replicator.tex}'s \texttt{DPROC} does) is deliberate: that
model needed \texttt{dstuck} versus \texttt{ddead} to route to two different
recoveries it fully derives; this model only needs to know \emph{whether}
a connected-but-misbehaving send happened at all, because the recovery
itself (\texttt{rhealing}) is the already-proven machinery, entered and left
in one step.

% ============================================================
\section{Initialization}
% ============================================================

\begin{schema}{Init}
  Link'
\where
  live' = \emptyset \\
  dirty' = \emptyset \\
  shown' = \emptyset \\
  linked' = clear \\
  reconnected' = clear \\
  wokenEarly' = clear \\
  reachable' = clear \\
  sendOk' = set \\
  regLock' = free \\
  rphase' = ridle \\
  discBackoff' = 0 \\
  wedgeBackoff' = 0 \\
  lInterval' = pdisc
\end{schema}

A single initial state: nothing live, nothing held, no display connected ---
\texttt{DISCONNECTED}, the quiet healthy idle Hub of \S4's own table ---
both counters at their base, the prober already at its relaxed pace, no
lock held.

% ============================================================
\section{Content Operations}
% ============================================================

An agent's push and withdrawal are content-owner actions on
\texttt{HubDisplay}.  Neither mentions \texttt{linked}, \texttt{reachable},
\texttt{regLock}, or \texttt{rphase}: no state of the display or of either
dialing actor can disable them, exactly the ``the Hub never controls the
display'' table of \S4 restated as an operation guard --- or rather, the
absence of one.

\begin{schema}{PushContent}
  \Delta Link \\
  s? : SCENE
\where
  live' = live \cup \{ s? \} \\
  dirty' = dirty \cup \{ s? \} \\
  shown' = shown \setminus \{ s? \} \\
  linked' = linked \\
  reconnected' = reconnected \\
  wokenEarly' = wokenEarly \\
  reachable' = reachable \\
  sendOk' = sendOk \\
  regLock' = regLock \\
  rphase' = rphase \\
  discBackoff' = discBackoff \\
  wedgeBackoff' = wedgeBackoff \\
  lInterval' = lInterval
\end{schema}

A scene that becomes live is marked owed and, if it had been believed
shown, that belief is withdrawn --- the new content has not been delivered
yet, whatever the connection state.

\begin{schema}{WithdrawContent}
  \Delta Link \\
  s? : SCENE
\where
  s? \in live \\
  live' = live \setminus \{ s? \} \\
  dirty' = dirty \setminus \{ s? \} \\
  shown' = shown \setminus \{ s? \} \\
  linked' = linked \\
  reconnected' = reconnected \\
  wokenEarly' = wokenEarly \\
  reachable' = reachable \\
  sendOk' = sendOk \\
  regLock' = regLock \\
  rphase' = rphase \\
  discBackoff' = discBackoff \\
  wedgeBackoff' = wedgeBackoff \\
  lInterval' = lInterval
\end{schema}

Disposing content before any display connects is the design's
\texttt{HELD}~$\to$~\texttt{DISCONNECTED} transition (\S4's table); nothing
in this operation depends on \texttt{linked}, so it applies identically
whether connected or not.

% ============================================================
\section{Environment}
% ============================================================

The display's own reachability and send-health change independently of the
Hub, exactly as \texttt{hub\_replicator.tex}'s \texttt{DisplayWedge} and
\texttt{DisplayCrash} model a display failing on its own.  These are the
\emph{only} operations that ever write a new value to \texttt{reachable} or
\texttt{sendOk} --- the textual fact I5 rests on.

\begin{schema}{GoReachable}
  \Delta Link
\where
  reachable = clear \\
  reachable' = set \\
  live' = live \\ dirty' = dirty \\ shown' = shown \\
  linked' = linked \\ reconnected' = reconnected \\
  wokenEarly' = wokenEarly \\ sendOk' = sendOk \\
  regLock' = regLock \\ rphase' = rphase \\
  discBackoff' = discBackoff \\ wedgeBackoff' = wedgeBackoff \\
  lInterval' = lInterval
\end{schema}

\begin{schema}{GoUnreachable}
  \Delta Link
\where
  reachable = set \\
  reachable' = clear \\
  live' = live \\ dirty' = dirty \\ shown' = shown \\
  linked' = linked \\ reconnected' = reconnected \\
  wokenEarly' = wokenEarly \\ sendOk' = sendOk \\
  regLock' = regLock \\ rphase' = rphase \\
  discBackoff' = discBackoff \\ wedgeBackoff' = wedgeBackoff \\
  lInterval' = lInterval
\end{schema}

\begin{schema}{GoSendOk}
  \Delta Link
\where
  sendOk = clear \\
  sendOk' = set \\
  live' = live \\ dirty' = dirty \\ shown' = shown \\
  linked' = linked \\ reconnected' = reconnected \\
  wokenEarly' = wokenEarly \\ reachable' = reachable \\
  regLock' = regLock \\ rphase' = rphase \\
  discBackoff' = discBackoff \\ wedgeBackoff' = wedgeBackoff \\
  lInterval' = lInterval
\end{schema}

\begin{schema}{GoSendFail}
  \Delta Link
\where
  sendOk = set \\
  sendOk' = clear \\
  live' = live \\ dirty' = dirty \\ shown' = shown \\
  linked' = linked \\ reconnected' = reconnected \\
  wokenEarly' = wokenEarly \\ reachable' = reachable \\
  regLock' = regLock \\ rphase' = rphase \\
  discBackoff' = discBackoff \\ wedgeBackoff' = wedgeBackoff \\
  lInterval' = lInterval
\end{schema}

% ============================================================
\section{The Replicator}
% ============================================================

The replicator dials only when it has something to deliver, exactly as the
real worker only calls \texttt{.get()} from inside a drained cycle.  A dial
that finds the display reachable either completes an ordinary already-
connected cycle or --- if the registry was not connected a moment ago ---
is the fresh connect that reconciles every held scene in the very same
step, discharging I2 by construction rather than by a separate pass.

\begin{schema}{RDialOk}
  \Delta Link \\
\where
  rphase = ridle \\
  dirty \neq \emptyset \\
  reachable = set \\
  linked' = set \\
  (linked = clear \implies
    (reconnected' = set \land dirty' = dirty \cup live)) \\
  (linked = set \implies
    (reconnected' = reconnected \land dirty' = dirty)) \\
  discBackoff' = 0 \\
  rphase' = rsending \\
  live' = live \\ shown' = shown \\ wokenEarly' = wokenEarly \\
  reachable' = reachable \\ sendOk' = sendOk \\ regLock' = regLock \\
  wedgeBackoff' = wedgeBackoff \\ lInterval' = lInterval
\end{schema}

A dial that finds nothing listening is exactly today's \texttt{RuntimeError}
from \texttt{ClientRegistry.get()}: no content is touched (nothing is lost,
\texttt{dirty} is untouched), \texttt{linked} drops, the disconnected
backoff advances toward its own cap, and the replicator moves to the wait.
This is the one and only operation that advances \texttt{discBackoff}.

\begin{schema}{RDialFail}
  \Delta Link
\where
  rphase = ridle \\
  dirty \neq \emptyset \\
  reachable = clear \\
  linked' = clear \\
  (discBackoff < 3 \implies discBackoff' = discBackoff + 1) \\
  (discBackoff = 3 \implies discBackoff' = discBackoff) \\
  rphase' = rwaitdisc \\
  regLock' = free \\
  live' = live \\ dirty' = dirty \\ shown' = shown \\
  reconnected' = reconnected \\ wokenEarly' = wokenEarly \\
  reachable' = reachable \\ sendOk' = sendOk \\
  wedgeBackoff' = wedgeBackoff \\ lInterval' = lInterval
\end{schema}

The cap is \texttt{3} here, standing in for the design's 120s ceiling;
\texttt{regLock' = free} is not a passthrough but the load-bearing fact of
\S6.3 stated at the one place it is decided --- entering the wait never
takes the lock.  \texttt{RDialFail}'s guard is disjoint from
\texttt{RDialOk}'s (\texttt{reachable} has exactly two values), so exactly
one of the two is ever the successor from \texttt{ridle} with work pending.

The wait is broken one of two ways.  \texttt{RWaitWoken} models
\texttt{ClientRegistry.wait\_for\_reconnect} returning \texttt{True}: the
event had been set by \emph{either} actor's successful dial, so the
replicator retries at once, without paying any more of the current delay.

\begin{schema}{RWaitWoken}
  \Delta Link
\where
  rphase = rwaitdisc \\
  reconnected = set \\
  rphase' = ridle \\
  reconnected' = clear \\
  wokenEarly' = set \\
  live' = live \\ dirty' = dirty \\ shown' = shown \\ linked' = linked \\
  reachable' = reachable \\ sendOk' = sendOk \\ regLock' = regLock \\
  discBackoff' = discBackoff \\ wedgeBackoff' = wedgeBackoff \\
  lInterval' = lInterval
\end{schema}

\texttt{RWaitTimeout} models the same call returning \texttt{False}: the
full current delay elapsed with no reconnect.  \texttt{wait\_for\_reconnect}
clears the event on \emph{either} exit, matching \texttt{Event.wait}
followed unconditionally by \texttt{.clear()} in the real code, so both
operations agree on that one point and differ only in \texttt{wokenEarly}.

\begin{schema}{RWaitTimeout}
  \Delta Link
\where
  rphase = rwaitdisc \\
  rphase' = ridle \\
  reconnected' = clear \\
  wokenEarly' = clear \\
  live' = live \\ dirty' = dirty \\ shown' = shown \\ linked' = linked \\
  reachable' = reachable \\ sendOk' = sendOk \\ regLock' = regLock \\
  discBackoff' = discBackoff \\ wedgeBackoff' = wedgeBackoff \\
  lInterval' = lInterval
\end{schema}

Once connected, a delivery either succeeds --- everything held is shown,
and the wedged-display backoff resets, since a clean cycle is unambiguous
in a way ``hasn't misbehaved again yet'' is not (\S6.3) --- or the send
itself fails against an already-open connection, which is the
\emph{different} condition this design is careful never to confuse with a
failed dial.

\begin{schema}{RSendCycle}
  \Delta Link
\where
  rphase = rsending \\
  sendOk = set \\
  dirty' = \emptyset \\
  shown' = shown \cup dirty \\
  wedgeBackoff' = 0 \\
  rphase' = ridle \\
  live' = live \\ linked' = linked \\ reconnected' = reconnected \\
  wokenEarly' = wokenEarly \\ reachable' = reachable \\ sendOk' = sendOk \\
  regLock' = regLock \\ discBackoff' = discBackoff \\
  lInterval' = lInterval
\end{schema}

\begin{schema}{RSendFail}
  \Delta Link
\where
  rphase = rsending \\
  sendOk = clear \\
  (wedgeBackoff < 1 \implies wedgeBackoff' = wedgeBackoff + 1) \\
  (wedgeBackoff = 1 \implies wedgeBackoff' = wedgeBackoff) \\
  rphase' = rhealing \\
  live' = live \\ dirty' = dirty \\ shown' = shown \\ linked' = linked \\
  reconnected' = reconnected \\ wokenEarly' = wokenEarly \\
  reachable' = reachable \\ sendOk' = sendOk \\ regLock' = regLock \\
  discBackoff' = discBackoff \\ lInterval' = lInterval
\end{schema}

The cap is \texttt{1} here, standing in for the design's 2.0s ceiling ---
deliberately much lower than \texttt{discBackoff}'s cap of \texttt{3}, the
concrete gap I4 is stated over.  \texttt{RSendFail} is the only operation
that advances \texttt{wedgeBackoff}, and it never touches
\texttt{discBackoff}; \texttt{RDialFail} is the only operation that
advances \texttt{discBackoff}, and it never touches \texttt{wedgeBackoff}.
That disjointness is I7, established by inspection exactly as
\texttt{hub\_replicator.tex}'s I1 is (``no mutator operation mentions
\texttt{clientLock}''): there is no third operation and no shared branch
that could touch both.

\texttt{RHeal} abstracts the reap/respawn/reconnect/re-mark recovery
\texttt{hub\_replicator.tex} already proves in full (its \texttt{Reap},
\texttt{Ensure}, \texttt{Remark}, \texttt{Reconn}): here it is one atomic
step, entered only from \texttt{rhealing} and always eventually leaving
it, touching nothing this spec cares about beyond returning control to the
cycle.

\begin{schema}{RHeal}
  \Delta Link
\where
  rphase = rhealing \\
  rphase' = ridle \\
  live' = live \\ dirty' = dirty \\ shown' = shown \\ linked' = linked \\
  reconnected' = reconnected \\ wokenEarly' = wokenEarly \\
  reachable' = reachable \\ sendOk' = sendOk \\ regLock' = regLock \\
  discBackoff' = discBackoff \\ wedgeBackoff' = wedgeBackoff \\
  lInterval' = lInterval
\end{schema}

% ============================================================
\section{The Liveness Prober}
% ============================================================

\texttt{DisplayLiveness} probes continuously and independently of the
replicator's own cycle --- it has no notion of ``dirty content,'' only of
whether the display answers.  Its dial goes through the same registry lock
the replicator's does, so its guard is exactly \texttt{regLock = free};
under the corrected design that guard is never actually contended, because
nothing ever holds the lock across a wait (\S Fidelity makes this concrete).
A successful probe can be the very dial that turns \texttt{linked} from
\texttt{clear} to \texttt{set} and reconciles held content --- ``ANY
successful \texttt{ClientRegistry.get()}'' in the design's own words, not
only the replicator's.

\begin{schema}{LDialOk}
  \Delta Link
\where
  regLock = free \\
  reachable = set \\
  linked' = set \\
  (linked = clear \implies
    (reconnected' = set \land dirty' = dirty \cup live)) \\
  (linked = set \implies
    (reconnected' = reconnected \land dirty' = dirty)) \\
  lInterval' = pconn \\
  live' = live \\ shown' = shown \\ wokenEarly' = wokenEarly \\
  reachable' = reachable \\ sendOk' = sendOk \\ regLock' = regLock \\
  rphase' = rphase \\
  discBackoff' = discBackoff \\ wedgeBackoff' = wedgeBackoff
\end{schema}

\begin{schema}{LDialFail}
  \Delta Link
\where
  regLock = free \\
  reachable = clear \\
  linked' = clear \\
  lInterval' = pdisc \\
  live' = live \\ dirty' = dirty \\ shown' = shown \\
  reconnected' = reconnected \\ wokenEarly' = wokenEarly \\
  reachable' = reachable \\ sendOk' = sendOk \\ regLock' = regLock \\
  rphase' = rphase \\
  discBackoff' = discBackoff \\ wedgeBackoff' = wedgeBackoff
\end{schema}

Neither operation is guarded by \texttt{rphase}: the prober runs regardless
of whatever cycle the replicator is in, which is the whole point of it
being a genuinely independent second actor rather than a second phase of
the same worker.

% ============================================================
\section{Invariants}
\label{sec:inv}
% ============================================================

None of the eight properties below is placed in \texttt{Link}'s own
predicate, for the same reason \texttt{hub\_replicator.tex} keeps its six
out of \texttt{Repl}: a predicate there would silently disable any
operation that would break it, masking exactly the defect the model exists
to catch.

\paragraph{I1 --- content never lost while disconnected.}
Every live scene is either owed or already believed delivered:
\[
  \forall s : SCENE \mid s \in live @ s \in dirty \lor s \in shown.
\]
No operation removes a scene from \texttt{dirty} without also adding it to
\texttt{shown} (\texttt{RSendCycle}) or removing it from \texttt{live} in
the same step (\texttt{WithdrawContent}); no operation ever adds to
\texttt{live} without adding the same scene to \texttt{dirty}
(\texttt{PushContent}).  The claim holds unconditionally, not only while
\texttt{linked = clear}: a scene mid-cycle, not yet delivered, is still
covered because it is still in \texttt{dirty}.

\paragraph{I2 --- reconcile exactly once.}
\texttt{dirty' = dirty \cup live} fires only inside \texttt{RDialOk} and
\texttt{LDialOk}, and only under the guard \texttt{linked = clear} ---
i.e.\ only on the step that turns the registry from disconnected to
connected.  Because \texttt{linked} is one shared boolean, at most one of
the two dialing operations can be the step that performs that transition
for a given disconnection; whichever one does not (because it finds
\texttt{linked} already \texttt{set}) takes the other branch, which leaves
\texttt{dirty} untouched by the reconcile.  The reconcile cannot be missed
either, because it is not a separate pass triggered \emph{after} the
transition: it is part of the same atomic step that performs the
transition.  A held scene believed already shown is the sharpest witness:
only the reconcile can put a scene into \texttt{dirty} without first
taking it out of \texttt{shown}, since every other producer of
\texttt{dirty} (\texttt{PushContent}) removes it from \texttt{shown} in the
same step.

\paragraph{I3 --- neither actor blocks or spins on a missing display.}
Every dial (\texttt{RDialOk}, \texttt{RDialFail}, \texttt{LDialOk},
\texttt{LDialFail}) is one Z operation and therefore resolves in one
bounded step by construction --- there is no schema in this model whose
enabling condition can remain true indefinitely without an intervening
state change, unlike a raw \texttt{time.sleep} which is not a state
transition at all.  \texttt{lInterval} is typed \texttt{SPEED}, so the
prober's pacing is always exactly one of the two named values; the
replicator's own interval between retries is governed by
\texttt{discBackoff} alone (\texttt{RWaitTimeout}'s only enabling condition
is \texttt{rphase = rwaitdisc}, entered only via \texttt{RDialFail}, which
is the only place \texttt{discBackoff} advances) --- there is no third,
unnamed pacing value either actor can fall into.

\paragraph{I4 --- the disconnected backoff reaches its own, higher, slow
steady state.}
\[
  \Diamond\,(discBackoff = 3)
\]
is reachable under sustained disconnection alone (no send ever attempted).
This is the specific defect this design fixes, restated as a reachability
property rather than an ``eventually bounded'' one: \S Fidelity shows that
collapsing the two counters makes this exact state unreachable, because
nothing then advances \texttt{discBackoff} past whatever the shared,
lower-capped counter tops out at.

\paragraph{I5 --- the Hub never issues a process-control call to the
display.}
No operation other than \texttt{GoReachable} ever writes \texttt{set} to
\texttt{reachable'} where \texttt{reachable = clear} beforehand ---
inspection of every schema above confirms each either leaves
\texttt{reachable} unchanged or is \texttt{GoReachable}/\texttt{GoUnreachable}
themselves, the two environment operations standing in for a human running
\texttt{lux display start}/\texttt{stop} or the OS service supervisor.  This
is the absolute, grep-checkable form of the operator's ruling (\S2 item 1
of the design), restated as the claim that the model's own carrier of
``a display becomes reachable'' has exactly one non-environment writer:
none.

\paragraph{I6 --- a display connect promptly breaks the replicator out.}
\[
  \Diamond\,(wokenEarly = set)
\]
is reachable: the replicator's wait can be broken by the shared event, not
only by the full delay elapsing.  \S Fidelity shows this is exactly the
property that a raw blocking sleep --- \texttt{RWaitWoken} deleted ---
makes unreachable, reproducing ``up-to-full-delay blindness'' precisely.

\paragraph{I7 --- the two backoffs never advance together.}
\texttt{RDialFail} is the only operation whose effect changes
\texttt{discBackoff}, and its frame leaves \texttt{wedgeBackoff} unchanged;
\texttt{RSendFail} is the only operation whose effect changes
\texttt{wedgeBackoff}, and its frame leaves \texttt{discBackoff} unchanged.
No operation instance can therefore advance both counters in one step ---
established by inspection, exactly as \texttt{hub\_replicator.tex}
establishes its own I1.  The mechanism is corroborated live (not simply
dead code) by the reachability of
\[
  \Diamond\,(discBackoff > 0 \land wedgeBackoff > 0)
\]
--- both counters independently nonzero, from two disjoint causes.

\paragraph{I8 --- deadlock-freedom; the wait never holds the registry
lock.}
\[
  \lnot\,\Diamond\,(rphase = rwaitdisc \land regLock = held).
\]
No operation in this specification ever sets \texttt{regLock' = held}: it
is \texttt{free} at \texttt{Init} and every subsequent schema either leaves
it unchanged or (\texttt{RDialFail}) explicitly re-asserts \texttt{free} at
the one place \S6.3's design decision is made.  The prober's own dial
(\texttt{LDialOk}/\texttt{LDialFail}) is therefore never disabled by the
replicator sitting in \texttt{rwaitdisc}, and the full state space is free
of deadlock (Verification, below).  \S Fidelity constructs the variant
where this fails, and shows the goal above becomes reachable and the
prober's dial becomes starved for as long as the replicator waits.

% ============================================================
\section{Verification}
\label{sec:verify}
% ============================================================

Type-checking is a \texttt{fuzz} pass:
\begin{verbatim}
  fuzz docs/display_linkage.tex
\end{verbatim}

I1, I5, and I8's lock claim are state properties, checked as reachability
of the negation (\emph{goal NOT found} means the invariant holds
everywhere reachable):
\begin{verbatim}
  # Invariant 1 (must report: goal NOT found)
  probcli docs/display_linkage.tex -model_check -nodead \
    -goal "card({s | s : SCENE & s : live &
           not(s : dirty) & not(s : shown)}) > 0" \
    -p DEFAULT_SETSIZE 2

  # Invariant 8, the lock claim (must report: goal NOT found)
  probcli docs/display_linkage.tex -model_check -nodead \
    -goal "rphase = rwaitdisc & regLock = held" \
    -p DEFAULT_SETSIZE 2
\end{verbatim}

I2, I3, I5, and I7 are established by inspection of the operation schemas
(Section~\ref{sec:inv}) and corroborated live by reachability goals that
must be \emph{found} --- the mechanism is exercised, not dead:
\begin{verbatim}
  # Invariant 2 corroboration (must report: goal FOUND)
  probcli docs/display_linkage.tex -model_check -nodead \
    -goal "card(dirty /\ shown) > 0" -p DEFAULT_SETSIZE 2

  # Invariant 3 corroboration -- both named prober speeds are live
  probcli docs/display_linkage.tex -model_check -nodead \
    -goal "lInterval = pdisc" -p DEFAULT_SETSIZE 2
  probcli docs/display_linkage.tex -model_check -nodead \
    -goal "lInterval = pconn" -p DEFAULT_SETSIZE 2

  # Invariant 7 corroboration (must report: goal FOUND)
  probcli docs/display_linkage.tex -model_check -nodead \
    -goal "discBackoff > 0 & wedgeBackoff > 0" -p DEFAULT_SETSIZE 2
\end{verbatim}

I4 and I6 are the two properties this design exists to fix, each checked as
reachability that must be \emph{found} in the corrected model:
\begin{verbatim}
  # Invariant 4 (must report: goal FOUND)
  probcli docs/display_linkage.tex -model_check -nodead \
    -goal "discBackoff = 3" -p DEFAULT_SETSIZE 2

  # Invariant 6 (must report: goal FOUND)
  probcli docs/display_linkage.tex -model_check -nodead \
    -goal "wokenEarly = set" -p DEFAULT_SETSIZE 2
\end{verbatim}

Deadlock freedom is the full-state-space check.  Every phase has an
always-available successor --- \texttt{PushContent}/\texttt{WithdrawContent}
require nothing about the display, \texttt{RWaitTimeout} requires nothing
beyond \texttt{rphase = rwaitdisc}, and each of \texttt{sendOk}'s two values
enables exactly one of \texttt{RSendCycle}/\texttt{RSendFail} --- so any
deadlock reported would be a genuine cyclic wait, not this model
terminating:
\begin{verbatim}
  # Deadlock freedom (must report: no deadlock / no counterexample)
  probcli docs/display_linkage.tex -model_check -p DEFAULT_SETSIZE 2
\end{verbatim}

% ============================================================
\section{Fidelity: reproducing the design's own defects}
\label{sec:fidelity}
% ============================================================

A model that cannot reproduce the defects it is meant to guard against
proves nothing.  This specification guards three; each is reproduced by a
small, named change to one operation, verified against a scratch copy of
this file with exactly that change applied and nothing else.

\paragraph{(a) Collapsing the tri-state outcome onto the shared backoff.}
The design's own root-cause bug (\S1, \S6.1): before the fix,
\emph{every} not-connected cycle advanced the same low-capped backoff a
wedged display uses, because the code could not tell ``never connected''
apart from ``connected but just healed.''  Reproduced here by changing
\texttt{RDialFail}'s effect from advancing \texttt{discBackoff} (cap 3) to
advancing \texttt{wedgeBackoff} (cap 1) instead --- the literal
``collapse onto the shared backoff'':
\begin{quote}
  \texttt{(wedgeBackoff < 1 \(\implies\) wedgeBackoff' = wedgeBackoff + 1)} \\
  \texttt{(wedgeBackoff = 1 \(\implies\) wedgeBackoff' = wedgeBackoff)} \\
  \texttt{discBackoff' = discBackoff}
\end{quote}
in place of the two \texttt{discBackoff} lines shown above.  The property
that separates the two designs is I4 itself:
\begin{verbatim}
  # must be FOUND for the correct spec, NOT FOUND for the collapsed variant
  probcli <spec>.tex -model_check -nodead \
    -goal "discBackoff = 3" -p DEFAULT_SETSIZE 2
\end{verbatim}
Run against a scratch copy with the one-line substitution above: the goal
returns \emph{not found} --- \texttt{discBackoff} is never touched by any
operation in that variant, so it never leaves \texttt{0}, while
\texttt{wedgeBackoff} climbs under sustained disconnection alone and
plateaus at its own cap of \texttt{1} --- the model's rendering of ``the
backoff plateaus at $\sim$2s and stays there forever,'' exactly the
symptom of \S1's confirmed root cause. Reachability of
\texttt{wedgeBackoff = 1} under a disconnection-only trace (no
\texttt{RSendFail} ever fired) is the positive confirmation of the same
collapse:
\begin{verbatim}
  # FOUND in the collapsed variant via RDialFail alone
  probcli <spec>.tex -model_check -nodead \
    -goal "wedgeBackoff = 1" -p DEFAULT_SETSIZE 2
\end{verbatim}

\paragraph{(b) A blocking sleep in place of the interruptible wait.}
The design's own second, load-bearing correction (\S6.3): a raw
\texttt{time.sleep(delay)} would commit the replicator to sleeping out the
\emph{entire} current delay --- up to 120s once backed off --- before
attempting another connect, regardless of how promptly the display actually
returns.  Reproduced here by deleting \texttt{RWaitWoken} from the
operation set entirely, leaving \texttt{RWaitTimeout} as the only way out
of \texttt{rwaitdisc}:
\begin{verbatim}
  # RWaitWoken removed; RWaitTimeout unchanged
\end{verbatim}
The property that separates the two designs is I6 itself:
\begin{verbatim}
  # must be FOUND for the correct spec, NOT FOUND for the blocking-sleep variant
  probcli <spec>.tex -model_check -nodead \
    -goal "wokenEarly = set" -p DEFAULT_SETSIZE 2
\end{verbatim}
With \texttt{RWaitWoken} removed, nothing in the model ever sets
\texttt{wokenEarly} to \texttt{set} --- the goal returns \emph{not found}.
Concretely: \texttt{reachable} can flip to \texttt{set}
(\texttt{GoReachable}) and the prober can dial successfully
(\texttt{LDialOk}, itself unaffected by the deletion) and set
\texttt{reconnected}, all while \texttt{rphase = rwaitdisc} --- but with no
\texttt{RWaitWoken} to consume that signal, the replicator remains parked
until \texttt{RWaitTimeout} eventually fires regardless, discovering the
display no sooner than a display that never reconnected at all would have
been retried.  This is the model's rendering of ``a display connect not
promptly breaking the replicator out'' --- exactly the second required
symptom, and exactly the failure mode \S6.3's own worked example describes
(``a user who runs \texttt{lux display start} right after the Hub has
climbed deep into its backoff would see held content stay blank for up to
two minutes'').

\paragraph{(c) Holding the registry lock across the wait (bonus control
for I8).} Not one of the two symptoms the mission requires, but the most
direct way to exhibit what \S6.3's own prose warns the lock-freedom
property prevents.  Reproduced by two changes: \texttt{RDialFail} sets
\texttt{regLock' = held} instead of \texttt{free}, and both
\texttt{RWaitWoken} and \texttt{RWaitTimeout} gain \texttt{regLock' = free}
in place of their passthrough. The distinguishing goal is I8's own:
\begin{verbatim}
  # must be NOT FOUND for the correct spec, FOUND for the lock-holding variant
  probcli <spec>.tex -model_check -nodead \
    -goal "rphase = rwaitdisc & regLock = held" -p DEFAULT_SETSIZE 2
\end{verbatim}
In the lock-holding variant this goal is reached immediately after the
first \texttt{RDialFail}, and for as long as the replicator remains
parked, \texttt{LDialOk}/\texttt{LDialFail}'s shared guard
\texttt{regLock = free} is false --- the prober is starved from dialling at
all, which in turn starves \texttt{reconnected} of the one signal that
could wake the replicator early, so the only escape becomes
\texttt{RWaitTimeout}: precisely the mechanism by which ``holding the lock
across the wait'' would defeat the very promptness I6 exists to guarantee,
stated here as a positive account of \emph{why} \S6.3 insists on
\texttt{wait\_for\_reconnect} not taking \texttt{ClientRegistry.\_lock},
not merely that it must not. This variant does not produce a literal ProB
deadlock --- \texttt{RWaitTimeout} remains an always-enabled escape, and
the content operations never consult \texttt{regLock} at all, so the state
space stays live -- it produces exactly the starvation the reachability
goal states, which is the honest and sufficient demonstration of the
claim; a spec that manufactured a true deadlock here would need to also
remove \texttt{RWaitTimeout}, which is a different, unrelated defect from
the one this control targets.

\end{document}
